\documentclass{situationist-preprint}

\title{The Contribution of the Observer}
\msid{AES-02-b}
\msfamily{AES-02}
\msdiscipline{Aesthetics}
\mstype{Companion manuscript}
\mskeywords{beauty; response-dependence; observer; secondary qualities; rigidification}

\begin{document}
\maketitle

\begin{abstract}
AES-02-a argues that beauty is a response-dependent property, located in the relation between object and observer. We ask what, precisely, the observer contributes. We distinguish three roles an observer might play: fixing which property beauty is, making a particular object beautiful, and detecting that it is. We argue that on the response-dependent account the observer, as a kind, plays the first role; the object alone plays the second; and individual observers play only the third. The distinction explains why the account is not subjectivism and why it does not make beauty depend on anyone's actually looking.
\end{abstract}

\section{Three roles}

Theorem~3.5 of AES-02-a states that an object is beautiful if and only if it is disposed to give aesthetic pleasure to qualified observers in suitable conditions, with qualified observers specified, in Definition~2.4, on Humean lines. The slogan that beauty lies in the relation between object and observer invites the question what each relatum contributes.

\begin{definition}\label{d:roles}
An observer may contribute to beauty in three ways: (a) \emph{individuation}, by fixing which property \enquote{beautiful} picks out; (b) \emph{instantiation}, by bringing it about that a given object has that property; (c) \emph{detection}, by registering that a given object has it.
\end{definition}

\begin{proposition}\label{p:roles}
On the account of AES-02-a, the kind of qualified observer contributes to individuation, the object's own features determine instantiation, and individual observers contribute only detection.
\end{proposition}

\begin{proof}
Individuation: by Definition~2.3 of AES-02-a, the property is fixed by a response, a kind of observer and a kind of condition; change the kind of observer and a different property is picked out. Instantiation: whether a given object has the disposition depends on its features, since a disposition is grounded in the categorical properties of its bearer, and by Principle~P3 of AES-02-a beauty supervenes on the object's non-aesthetic features. No observer's response brings the disposition into being. Detection: an individual observer's pleasure, when she is qualified and conditions are suitable, manifests the disposition and is evidence of it; when she is not, her pleasure is evidence of nothing.
\end{proof}

The proposition locates the truth in the saying that beauty is in the eye of the beholder. The beholder, taken as a kind, is in the property, in the way normal human vision is in the property of being red \citep[II.viii]{Locke1690}. But no particular beholder is in it, and that is why subjectivism, which gives the individual observer the role of instantiation, fails the condition of correctability (Lemma~3.1 of AES-02-a).

\section{Two consequences}

First, the account does not make beauty depend on anyone's looking. An object unobserved has the disposition if it has the grounding features, as an unseen tomato is red. Second, if the account is rigidified, so that the observers in question are qualified observers as they actually are \citep{Johnston1989}, then a change in human sensibility would change what people enjoy and not what is beautiful. Rigidification is not an ad hoc repair: it reflects the fact that the kind of observer contributes to individuation, and individuation is fixed by the observers we have, not by those we might become.

\begin{objection}[Projection]
If the kind of observer fixes the property, beauty is still our projection onto things, not a feature they have.
\end{objection}

\begin{reply}
Fixing which property a word picks out by reference to a response is not projecting a response onto objects. The disposition picked out is had by objects independently of anyone's projecting, and judgments ascribing it can be false. \citet{McDowell1985} made the parallel point for secondary qualities: a property can be essentially such as to be experienced in a certain way and still be there to be experienced.
\end{reply}

\section{Conclusion}

The observer contributes to beauty as a kind, by fixing which property is in question, and as an individual only as a detector. The object determines whether it is beautiful. This is the precise sense in which beauty is located in the relation between object and observer, and it is why the relation leaves beauty out of the eye of any particular beholder.

\begin{thebibliography}{99}\small
\bibitem[Johnston(1989)]{Johnston1989} Johnston, M. (1989). Dispositional theories of value. \emph{Proceedings of the Aristotelian Society, Supplementary Volume}, 63, 139--174.
\bibitem[Locke(1690)]{Locke1690} Locke, J. (1690). \emph{An Essay Concerning Human Understanding}. Ed. P.~H. Nidditch. Clarendon Press, 1975.
\bibitem[McDowell(1985)]{McDowell1985} McDowell, J. (1985). Values and secondary qualities. In T. Honderich (Ed.), \emph{Morality and Objectivity} (pp.~110--129). Routledge \& Kegan Paul.
\end{thebibliography}

\end{document}
